\section{Method}
\label{sec:method}

\paragraph{Layer-wise probing.}
For a frozen ViT-L encoder ($L{=}24$ blocks, $D{=}1024$, $H{=}16$ heads), we extract
representations at 25 positions ($\ell{=}0$: patch embedding; $\ell{=}1\ldots24$: block outputs).
DINOv2 uses the CLS token; JEPA models use mean-pooled patch tokens.
A one-vs-rest logistic regression ($C{=}1$, L-BFGS) is trained per layer to predict each of the
7 binary affordance labels; mAP is reported on the held-out test set.

\paragraph{Attention head ablation.}
At peak layer $\ell^*$, we ablate head $h$ by zeroing its $d_h{=}64$ output columns before
the output projection~\cite{michel2019sixteen}; importance is
$\Delta_h = \text{mAP}_\text{full} - \text{mAP}_{-h}$.
We sweep top-$k$ head subsets to quantify the efficiency--accuracy tradeoff; concentration is measured as each head's share of the total positive $\Delta_h$ mass.

\paragraph{Experimental setup.}
UMD Part Affordance Dataset~\cite{myers2015affordance} (7 affordance classes): stratified 2,000-image subsample
(80/10/10 split, $224{\times}224$, ImageNet normalization).
Static images are replicated $T{=}8$ times as pseudo-video for JEPA models.
A random ViT-L establishes the chance-level floor.
All encoders are \textbf{frozen}; only probe weights are trained.
